\begin{frame}[t]{\ev{\tagA\,\tagP}\surf{COUNT ONCE}Each fan-in operation needs its own check.}
\vspace{.1cm}
\begin{center}
\begin{tikzpicture}[x=1cm,y=1cm,every node/.style={font=\fontsize{11.5}{14}\selectfont},
  op/.style={state,anchor=west,text width=1.8cm,minimum height=.74cm,font=\fontsize{11.5}{14}\selectfont\bfseries},
  ret/.style={candidate,anchor=west,text width=4.5cm,minimum height=.74cm},
  chk/.style={evidence,anchor=west,text width=5.6cm,minimum height=.74cm},
  head/.style={anchor=west,font=\fontsize{9}{10}\selectfont\bfseries,text=Teal}]
\node[head] at (-.05,.75) {FAN-IN};
\node[head] at (2.75,.75) {RETURNS};
\node[head] at (8.15,.75) {CHECK};
\foreach \y/\o/\r/\c in {0/Select/one candidate/{the chosen one, on held-out data},-.95/Compose/a merged artifact/{the merged version itself},-1.9/Pool/one pooled estimate/{noise and shared errors}}{
  \node[op] (o) at (0,\y) {\o};
  \node[ret] (r) at (2.75,\y) {\r};
  \node[chk] (c) at (8.15,\y) {\c};
  \draw[flow] (o)--(r);\draw[flow] (r)--(c);
}
\end{tikzpicture}
\end{center}
\vspace{.3cm}
{\centering\fontsize{11.5}{14}\selectfont\textbf{Why test the merge?}\enspace Each change starts a 1\,GB worker; the server has 2\,GB.\par}
\begin{center}
\begin{tikzpicture}[x=1cm,y=1cm,every node/.style={font=\fontsize{11.5}{14}\selectfont,minimum height=.6cm,align=center}]
\node[evidence,text width=1.75cm] at (0,0) {A+B \yes};
\node[evidence,text width=1.75cm] at (2.4,0) {A+C \yes};
\node[evidence,text width=1.75cm] at (4.8,0) {B+C \yes};
\node[bad,text width=3.3cm] at (8.0,0) {A+B+C\,=\,3\,GB \no};
\end{tikzpicture}
\end{center}
\sources{Illustrative capacity example, not a recorded incident. Union and intersection are further fan-in rules (Q\&A).}

\qadetail{The three rows describe decision semantics, not three implemented end-to-end systems. Selection needs a fixed evaluator over eligible candidates; argmax returns the candidate while max returns only its score, so keep candidate identity and a deterministic tie rule. Selection is adaptive, so the chosen candidate needs fresh evaluation for generalization; a deterministic regression check can still establish its asserted behavior. Coverage is not selected success (the patch-selector point from the software-engineering example). Code merge is a partial operation on edits from the same base: it may be undefined for conflicting edits or depend on conflict resolution, and it is checked on its final output. No universal associative or commutative code-merge law is claimed. The 1 GB example is illustrative and exact, not a recorded software incident: with U=1 for pass, Gamma_ij=U({i,j})-U({i})-U({j})+U(empty) gives zero pair interactions and a failing triple. Pooling requires one common target, calibrated information and a dependence model; it cannot combine different artifacts simply because each returns a number. The full menu also has union for disjoint shards (check coverage and duplicates), intersection for common findings (agreement can reflect a shared error) and a weighted mean with weights summing to one. ReLU is a pointwise transform, not a fan-in rule.}
\note{[1:10] Once trials return receipts, we combine them. That step is fan-in, and it is not one operation. If the trials tried different ideas, we select one candidate; that is the patch selector from the software example. Selection is itself a search, so we check the winner again, on held-out data. If the trials edited different parts of a program, we compose their edits, and the merged version needs its own test. A made-up example: each of three changes starts a one-gigabyte worker, and the server has two gigabytes. Each change alone fits. Every pair fits. All three together need three gigabytes and fail. No pair test covered that state. Finally, if the trials estimated the same quantity, we pool them, and the weights must account for errors they share. Next: what that sharing costs.}
\end{frame}
